% chapters/15_inference.tex
% Part of: AI / LLM / Agents / Hardware Notes

\section{Inference \& Decoding Strategies}

\subsection{Greedy, beam search, sampling}

At each decoding step the model produces a probability distribution over the vocabulary.
The decoding strategy determines how a token is selected from that distribution.

\textbf{Greedy decoding} always selects $\arg\max_t P(t \mid \text{context})$. It is
fast and deterministic but tends to produce repetitive, low-entropy output because the
model commits to the single most probable path at every step, discarding all alternatives.

\textbf{Beam search} maintains the top-$B$ candidate sequences (``beams'') in parallel.
At each step every beam is extended with its top-$B$ continuations, and the global
top-$B$ are kept. Final output is the highest-scoring complete sequence. Beam search
trades $O(B)$ memory for more globally coherent text and is the default for machine
translation; it tends to produce overly conservative text in open-ended generation.

\textbf{Sampling} draws from the full (or truncated) distribution, enabling varied and
creative output at the cost of occasional incoherence. Temperature, top-$k$, and top-$p$
control where sampling takes place (see \S below).

\begin{intuitionbox}
  Greedy $\approx$ beam search with $B=1$; both are deterministic.
  Sampling with temperature $T \to 0$ converges to greedy; $T \to \infty$ becomes uniform.
\end{intuitionbox}

\subsection{Temperature, top-k, top-p}

These three hyperparameters reshape the distribution before sampling.

\textbf{Temperature} $T$ scales the logits before the softmax:
\[
  P'(t) = \frac{\exp(z_t / T)}{\sum_{t'}\exp(z_{t'} / T)}
\]
$T < 1$ sharpens the distribution (more deterministic); $T > 1$ flattens it (more
exploratory). Typical creative-writing values: $T \approx 0.7$--$1.0$; coding: $T \approx
0.2$.

\textbf{Top-$k$ sampling} keeps only the $k$ highest-probability tokens, renormalises,
and samples. Prevents extremely unlikely tokens from being selected but the hard cutoff
can still include incoherent tokens if the distribution is flat.

\textbf{Top-$p$ (nucleus) sampling} (Holtzman et al.\ 2020) keeps the smallest set of
tokens whose cumulative probability $\geq p$, then samples. The effective vocabulary size
adapts to the confidence of the distribution: when the model is certain, only a few tokens
make the nucleus; when uncertain, many tokens share probability mass and a larger nucleus
is used. $p = 0.9$--$0.95$ is a common default.

\begin{notebox}
  Top-$p$ is generally preferred over top-$k$ for open-ended generation because it
  adapts automatically. Using both simultaneously ($p$ + $k$) is also common as an
  additional safety net.
\end{notebox}

\subsection{KV cache mechanics}

During autoregressive generation, computing attention for every new token would require
re-processing all preceding tokens. The \textbf{key-value (KV) cache} avoids this.

\textbf{How it works.} On the first forward pass (prefill phase), the model computes
$K$ and $V$ matrices for all input tokens and caches them. For every subsequent token
generated (decode phase), only the new token's $Q$, $K$, $V$ are computed; the new
$K$, $V$ vectors are appended to the cache and the full attention is computed using the
stored history. No recalculation of past keys and values is needed.

\textbf{Memory cost.} KV cache size $\propto$ \texttt{num\_layers} $\times$
\texttt{num\_heads} $\times$ \texttt{head\_dim} $\times$ \texttt{seq\_len} $\times$
\texttt{batch\_size}. For a 70B-parameter model serving 32K-token sequences in a batch
of 16, KV cache can easily exceed 30\,GB --- the dominant memory consumer at inference
time, not the weights themselves.

\textbf{KV cache with GQA/MQA.} Grouped-Query Attention (GQA) and Multi-Query Attention
(MQA) reduce the number of distinct $K$/$V$ heads, directly shrinking the KV cache
without changing the model's output quality much. This is why Llama\,3 and Gemma adopt GQA.

\textbf{PagedAttention and vLLM.} Traditional KV cache is allocated contiguously in GPU
memory, causing two problems: (1) fragmentation --- reserved but unused space; (2) no
sharing between sequences that share a common prefix (e.g.\ system-prompt batches).

\begin{paperbox}[title={PagedAttention (Kwon et al.\ 2023)}]
  Borrows OS virtual memory ideas. The KV cache is divided into fixed-size
  \emph{blocks} (``pages'') that are mapped via a block table, analogous to page
  tables in virtual memory. Blocks are allocated on demand and need not be contiguous
  in physical GPU memory. Copy-on-write enables prefix sharing: multiple sequences
  sharing the same system prompt use the same physical cache pages until they diverge.
  \textbf{vLLM} is the production serving library built on PagedAttention.
\end{paperbox}

PagedAttention reduces memory fragmentation to under 4\,\% and roughly doubles serving
throughput compared with static pre-allocation.

\subsection{Speculative decoding}

\textbf{Bottleneck insight.} LLM decode speed is memory-bandwidth-bound, not
compute-bound. Each decode step loads the full model weights from HBM to feed just a
single token --- utilisation of GPU tensor cores is low. Speculative decoding reframes
the problem to generate multiple tokens per weight-load.

\textbf{Algorithm.}
\begin{enumerate}
  \item A small \emph{draft model} (e.g.\ 7B) autoregressively generates $k$ candidate
    tokens cheaply.
  \item The large \emph{target model} (e.g.\ 70B) processes all $k$ tokens in a single
    parallel forward pass, producing probability distributions for each position.
  \item Each draft token is accepted or rejected via token-level rejection sampling: if
    the draft distribution $q$ is ``consistent'' with the target $p$, the token is
    accepted; otherwise it is replaced with a corrected sample from $p$.
  \item The process is unbiased: the accepted sequence has the exact same distribution
    as if the target model had generated it greedily/sampled independently.
\end{enumerate}

Typical speedup: $2$--$3\times$ on tasks where the draft model is accurate (code
generation, continuation). On highly creative tasks with low draft acceptance, gains are
minimal.

\begin{paperbox}[title={Multi-token Prediction (Gloeckle et al.\ 2024)}]
  Trains the model to predict the next $n$ tokens simultaneously using independent
  heads. Unlike speculative decoding (which uses a separate draft model), multi-token
  prediction is baked into the architecture. It functions as a form of
  \emph{self-speculation}: the base model is its own draft. It also improves
  sample-efficiency during training by providing $n$ gradient signals per forward pass.
\end{paperbox}

\subsection{Continuous batching}

\textbf{Static batching} groups requests into a fixed batch and runs them together.
When some sequences finish early and others continue, GPU cycles are wasted waiting for
the longest sequence to complete before swapping in new work.

\textbf{Continuous batching} (also called \emph{iteration-level scheduling} or
\emph{in-flight batching}) solves this by decoupling batch composition from individual
decode steps. After every iteration, completed sequences are evicted from the batch and
new requests are inserted immediately, filling the vacant slots.

\begin{intuitionbox}
  Continuous batching treats the GPU like a highway with on-ramps: new cars (requests)
  can join mid-journey as soon as a lane (batch slot) opens up, rather than waiting for
  all cars to exit at the same time.
\end{intuitionbox}

Combined with PagedAttention (variable-length KV cache pages), continuous batching
enables high-throughput serving at arbitrary sequence lengths. vLLM, TensorRT-LLM, and
SGLang all implement continuous batching as their default scheduling strategy.

\subsection{Real-time action chunking for robot deployment}

\begin{intuitionbox}
VLA models deployed on real robots face a latency problem. The model must predict a sequence of actions (a ``chunk'') fast enough that the robot does not stall between chunks. A naïve chunk boundary causes a visible pause and a sharp joint-angle jump. This is distinct from the LLM serving problem but is fundamentally an \emph{inference scheduling} challenge.
\end{intuitionbox}

\subsubsection{The action-chunking gap}

A VLA predicts an action chunk $[a_0, \ldots, a_H]$ covering a horizon $H$ timesteps. While the robot executes this chunk, the model must begin predicting the \emph{next} chunk. Because the model has inference latency $\Delta t_\text{inf}$, the second chunk is ready only at time $\Delta t_\text{inf}$ after the query was sent. If the robot has already consumed $k$ actions from the first chunk, a naïve switch from $a_k$ to $a_0'$ (the start of the second chunk) causes a discontinuous jump in joint angles — visible as a jitter or stall.

The problem is compounded by the \textbf{overlap region}: when predicting the second chunk, the model does not know which sub-sequence of the first chunk the robot is currently executing, so it cannot smoothly continue from the right state.

\subsubsection{Training-time action conditioning (TTAC)}

Proposed in the paper \emph{Training-Time Action Conditioning for Efficient Real-Time Chunking} and adopted in the Size Zero humanoid system. The core idea is to treat the overlap as an \textbf{imputing} problem:

\begin{enumerate}
  \item During \textbf{training}, randomly mask a prefix of the target action sequence and condition the model on the \emph{known} suffix of the previous chunk (the actions already being executed). The model learns to predict only the \emph{remaining} unknown portion, consistent with the known prefix.
  \item A \textbf{simulated inference delay} $\hat{\Delta t}$ is injected during training so the model learns to infer the missing overlap under realistic latency conditions.
  \item At \textbf{inference time}, the known suffix of the currently-executing chunk is passed as conditioning, and the model fills in only the remaining unseen actions — ensuring the predicted trajectory smoothly continues from the robot's current execution state.
\end{enumerate}

\begin{paperbox}[title={Training-Time Action Conditioning (TTAC)}]
Key insight: action chunking discontinuity is not a hardware problem but a \emph{conditioning} problem. By training the model to condition on partially-known action prefixes, the overlap region becomes a smooth interpolation rather than a hard boundary. No change to the robot hardware or inference pipeline is required — only a modified training objective.
\end{paperbox}

\subsubsection{Client-server architecture for real-time control}

For humanoid whole-body control the system is split into two processes:

\begin{itemize}
  \item \textbf{Client} (on the robot): captures observations (RGB, joint states, tactile, IMU) at 30 Hz and streams them to the server. Consumes the latest available action from a shared action buffer and sends joint commands to the low-level controller.
  \item \textbf{Server} (GPU desktop): runs two concurrent loops:
    \begin{itemize}
      \item \textbf{Control loop}: continuously reads the latest action from the inference buffer and forwards it to the client — always at 30 Hz regardless of inference speed.
      \item \textbf{Inference loop}: runs TTAC-conditioned VLA inference to produce the next action chunk, using knowledge of currently-executing actions to guarantee trajectory continuity.
    \end{itemize}
\end{itemize}

This architecture decouples the \emph{control rate} (fixed, robot-determined) from the \emph{inference rate} (GPU-bound), so inference latency never causes visible pauses or joint-angle jumps.

\begin{notebox}
\textbf{Observed effect}: without TTAC, the joint-angle trajectory shows a sharp discontinuity at every chunk boundary. With TTAC, the predicted trajectory's start smoothly matches the end of the previous trajectory, eliminating jitter. The control loop ensures the robot always receives a valid command even while inference is in-flight.
\end{notebox}
